\begin{abstract}
薄膜材料实验需要协调溶液制备、液体运输、涂覆与后处理，可靠自动化不仅要求机器人完成操作，还要求最终样品质量合格。本文提出以薄膜制备为主要应用的混合机器人框架，由异构工作站机械臂与移动运输平台协同执行材料实验。约束感知经典规划用于站间运输，姿态感知强化学习用于器材抓取，结合域随机化与少量真实示范正则化微调的行为克隆用于精密液体操作。具备失败恢复机制的分层状态机通过明确的完成、超时与恢复条件编排可复用技能。带时间戳的虚实状态同步作为监测与轨迹比较的辅助接口，不构成预测闭环控制。已有独立技能评估中，域随机化结合少样本微调取得82.7\%的合并回合成功率，该数值不代表完整薄膜制备成功率。薄膜实验已完成，但人工、固定轨迹与学习控制在膜厚均匀性、接触角、粗糙度及耐腐蚀性能上的定量对比仍需补充。本文进一步明确未见条件迁移测试及完整流程质量验收方案，并将梯度稀释定位为液体操作技能库的扩展应用。
\end{abstract}
